\chapter{Evaluation: Intent-Aware Path Selection}
\label{ch:p1eval}

This chapter evaluates the Part~I selectors on held-out pairs, hours, and intents. We first fix the
experimental setup and metrics (\cref{sec:p1eval:setup}). We then answer the four questions the design
raises, in turn: does intent conditioning need modulation rather than concatenation
(\cref{sec:p1eval:ablation})? Does the conditioned policy genuinely re-rank paths as the intent changes
(\cref{sec:p1eval:intent})? Does it retain the near-oracle-quality-at-low-probing property of the
selector it builds on (\cref{sec:p1eval:overhead})? And where does even the best single-path choice stop
being enough (\cref{sec:p1eval:ceiling})? The last result is the bridge to Part~II.

\section{Experimental Setup}
\label{sec:p1eval:setup}

\paragraph{Environment.} Experiments use one BRITE-generated SCION topology of 90~ASes and 528~inter-AS
links, partitioned into 2~ISDs with 6~core ASes. Beaconing (\cref{sec:p1impl:pipeline}) discovers
1004~core segments and 1127~up/down segments, which compose into candidate paths for 8010~routable
source--destination pairs, each with up to 30~paths. The network is driven by the 28-day hourly demand
model with per-link aggregation (\cref{sec:p1design:env}).

\paragraph{Train/test split.} Training uses the first 14~days; all numbers in this chapter are measured
on the last 14~days (336~held-out hours), a window the agents never trained on. Unless noted, each cell
of a results table aggregates 32~held-out pairs across those 336~hours, i.e.\ 10752~independent decision
contexts.

\paragraph{Intents.} The conditional agent is trained across the reward profiles of
\cref{sec:p1design:film} and evaluated on the four \emph{selection-relevant} distinctive profiles:
throughput (\emph{Throughput}), latency-averse (\emph{Low-Latency}), loss-averse (\emph{Low-Loss}), and
a balanced extreme (\emph{Balanced}). The probe-frugal profiles are excluded from the conditioning
studies because $w_{\text{probe}}$ is path-independent for a single-path selector and so cannot change
the chosen path; probing cost is instead the subject of \cref{sec:p1eval:overhead}.

\paragraph{Metrics.} For each method we report the intent-weighted reward of \cref{eq:p1reward} and the
achieved per-flow goodput and latency. Intent alignment is assessed two ways: by reward under each
intent, and by \emph{behavioral divergence} -- the degree to which the chosen path changes as the intent
vector changes, holding the decision context fixed. Concretely, for each context we count the distinct
paths a method selects across the four intents and normalize to $[0,1]$, so a method that ignores intent
scores~0 and one that picks a different path for every intent scores~1. Probing cost is measured by the
per-selection probe accounting of \cref{sec:p1impl:baselines}.

\section{Does Conditioning Need Modulation? The FiLM Ablation}
\label{sec:p1eval:ablation}

The central design claim of Part~I is that intent must \emph{modulate} the per-path features, not merely
be concatenated to the state (\cref{sec:p1design:concat}). We test it by walking up the agent ladder of
\cref{sec:p1impl:agents}: a Flat DQN with a single baked-in objective, a Conditional-Concat agent that
receives the intent in its value stream, and the Conditional-FiLM agent that modulates the advantage
stream with it. \Cref{tab:p1eval:ablation} reports the mean reward each agent attains under each intent,
together with its adaptivity (mean behavioral divergence across intents).

\begin{table}[t]
  \centering
  \caption{Intent-conditioning ablation. Mean intent-weighted reward under each of the four intents, and
  \emph{Adaptivity} -- the mean behavioral divergence across intents (0~=~intent-invariant path choice,
  1~=~a different path per intent). Each reward cell aggregates 10752~held-out decisions.}
  \label{tab:p1eval:ablation}
  \begin{tabular}{lccccc}
    \toprule
    Method & Throughput & Low-Latency & Low-Loss & Balanced & Adaptivity \\
    \midrule
    Flat DQN            & 0.333 & 0.701 & 0.868 & 0.561 & 0.011 \\
    Conditional-Concat  & 0.979 & 0.708 & 0.937 & 0.894 & 0.050 \\
    Conditional-FiLM    & 0.980 & 0.732 & 0.939 & 0.894 & 0.128 \\
    \bottomrule
  \end{tabular}
\end{table}

Three findings stand out. First, the unconditioned Flat DQN collapses onto a single trade-off: it scores
well on Low-Loss (0.868) but poorly on Throughput (0.333) and Balanced (0.561), and its adaptivity of
0.011 confirms it selects essentially the same path regardless of the intent it is asked to serve -- as
it must, having no intent input. Second, both conditional agents recover strong reward on every intent,
because both receive the intent vector. The decisive difference is \emph{how much they act on it}: the
FiLM head re-ranks paths far more than concatenation, with an adaptivity of 0.128 against 0.050 -- a
factor of roughly 2.5 -- and a correspondingly higher choice entropy (0.31 versus 0.13~bits). This is the
predicted signature of common-mode cancellation: an intent that enters only the value stream shifts all
path scores together and largely survives the $\arg\max$ unchanged, so Concat under-reacts even though it
technically ``takes intent as input.''

Third, and most telling, the extra re-ranking that FiLM performs is not idle motion: it converts into
reward exactly on the intent that requires re-ranking. On Throughput, Low-Loss, and Balanced the two
conditional agents tie, because those objectives are served by the same high-bandwidth paths the value
estimate already favors. On Low-Latency -- the intent that must steer \emph{away} from the high-bandwidth
default toward short, lightly loaded paths -- FiLM leads Concat, 0.732 against 0.708. We find that FiLM
matches or exceeds concatenation on every intent while never falling below it, and that the behavioral
divergence exposes the mechanism gap even where the reward gap is small.

\section{Intent Alignment}
\label{sec:p1eval:intent}

Adaptivity establishes that the FiLM policy changes its choice with intent; this section shows that it
changes it in the \emph{right} direction. We evaluate the single trained FiLM checkpoint under each
intent and measure both the reward it earns and the metrics of the paths it selects.

\begin{figure}[t]
  \centering
  \includegraphics[width=0.72\linewidth]{p1eval_intent_heatmap.png}
  \caption{Reward matrix $R(\text{intent told}, \text{objective scored})$ for the FiLM policy. Cell
  color is the reward relative to its column mean, so a red diagonal means that conditioning on an intent
  is the best way to score well under that same objective. The diagonal dominates the two objectives that
  trade off against the high-bandwidth default (Throughput, Low-Latency) and ties on the objectives that
  do not (Low-Loss, Balanced).}
  \label{fig:p1eval:heatmap}
\end{figure}

\Cref{fig:p1eval:heatmap} reports the reward $R(i, j)$ obtained when the agent is told intent~$i$ but its
choice is scored under objective~$j$; each cell is colored by its deviation from the column mean. Reading
down a column asks: to do well under this objective, which intent should the agent be given? For
Throughput the answer is unambiguous -- being told ``Throughput'' scores 0.98, far above any mismatched
intent, and telling the agent the wrong intent (Low-Latency) drops it to 0.75. For Low-Latency the
matching intent again wins its column (0.73), and the mismatch penalty is largest here, precisely because
latency is the objective that most demands re-ranking. For Low-Loss and Balanced the diagonal is a
near-tie: low-loss, low-delay paths coincide with the high-bandwidth paths that most intents already
prefer, so several intents score alike. The conditioning helps most, in other words, exactly where the
objectives genuinely conflict.

\begin{figure}[t]
  \centering
  \includegraphics[width=\linewidth]{p1eval_intent_boxplots.png}
  \caption{Distributions of the chosen path's metrics as the FiLM policy is conditioned on each intent,
  over a fixed set of decision contexts. The Low-Latency intent lowers the chosen latency, the Throughput
  intent raises the chosen goodput, and the Low-Loss intent raises the chosen path's trust.}
  \label{fig:p1eval:boxplots}
\end{figure}

\Cref{fig:p1eval:boxplots} shows the same behavior in the physical metrics of the chosen paths. Holding
the decision contexts fixed and varying only the intent, the Low-Latency intent shifts the chosen-latency
distribution down (median near \SI{12.5}{\milli\second} against roughly \SI{17}{\milli\second} for the
other intents), the Throughput intent shifts the chosen-goodput distribution up, and the Low-Loss intent
raises the trust of the chosen path. A single set of weights, conditioned on different intents, re-ranks
the same candidate paths along whichever axis the intent emphasizes -- the modulative behavior the FiLM
head was designed to produce.

\section{Probing Overhead}
\label{sec:p1eval:overhead}

A path selector is only useful if its selection quality does not come at the price of measuring every
path on every decision. Using the probing accounting of \cref{sec:p1impl:baselines}, we compare the FiLM
selector against heuristics that must probe exhaustively, under the Balanced intent so that goodput and
reward are directly comparable.

\begin{figure}[t]
  \centering
  \includegraphics[width=0.72\linewidth]{p1eval_probing.png}
  \caption{Achieved goodput against probing cost per selection (top-left is better). The FiLM selector
  reaches the goodput of exhaustive widest-path selection while probing a single path instead of all~30,
  cutting probing cost by roughly \(40\times\). Latency-only heuristics probe cheaply but select poorly.}
  \label{fig:p1eval:probing}
\end{figure}

\Cref{fig:p1eval:probing} places each method in the quality--cost plane. Widest-path attains the highest
goodput (\SI{8.24}{\giga\bit\per\second}) but must issue a full bandwidth probe on all 30~candidate
paths, at \SI{5388}{\milli\second} of probing per selection. The FiLM selector reaches
\SI{8.15}{\giga\bit\per\second} -- within about 1\% of widest-path -- while inspecting only the single
path it selects, at \SI{135}{\milli\second} per selection: a reduction in probing cost of roughly
$40\times$ at matched goodput. Because it probes so little, its intent-weighted reward (0.894) even edges
out widest-path (0.891), which pays the probing penalty 30~times over. The single-metric heuristics
probe cheaply but choose badly: shortest-path and lowest-latency, which never inspect bandwidth, deliver
only \SI{5.8}{\giga\bit\per\second} and \SI{6.1}{\giga\bit\per\second} respectively, roughly 40\% less
goodput than FiLM. We thus reproduce, for the conditioned selector, the near-oracle-quality-at-low-probing
result of Bourak~et~al.~\cite{Bourak2025}.

\section{The Single-Path Performance Ceiling}
\label{sec:p1eval:ceiling}

The results so far show that one conditioned policy can serve the full intent spectrum at low probing
cost. This final experiment asks the question that motivates Part~II: how far does the best single-path
choice carry an application as the network fills up? We bin every held-out decision by its realized
congestion -- the mean utilization across the flow's candidate paths at that hour -- and plot the goodput
each method achieves against it (\cref{fig:p1eval:ceiling}).

\begin{figure}[t]
  \centering
  \includegraphics[width=0.72\linewidth]{p1eval_ceiling.png}
  \caption{Achieved goodput against realized congestion (mean path utilization). The FiLM selector tracks
  the widest-path ceiling at every load, but that ceiling itself falls by about 25\% from light to heavy
  congestion: past a point, no single path can carry the demand. Load-blind heuristics degrade far more
  steeply.}
  \label{fig:p1eval:ceiling}
\end{figure}

Two effects appear. First, the FiLM selector tracks the widest-path ceiling across the entire congestion
range, confirming that it makes near-optimal use of whatever single-path capacity exists, while the
load-blind heuristics fall away sharply -- shortest-path and lowest-latency lose about 58\% of their
goodput from the lightest to the heaviest congestion bin, and random collapses entirely. Second, and
crucially, \emph{the ceiling itself descends}: even FiLM and widest-path drop from about
\SI{9.1}{\giga\bit\per\second} at light load to about \SI{6.8}{\giga\bit\per\second} at heavy load, a
loss of roughly 25\%. This is the single-path performance ceiling. When the shared bottleneck is
congested, the best available path is still one path, and one path can carry only so much; no amount of
intent-aware selection recovers the capacity that a single link cannot provide. Escaping this ceiling
requires spreading a flow across several paths at once -- the problem Part~II takes up.

\section{Summary}
\label{sec:p1eval:summary}

The Part~I evaluation supports four claims. One conditioned scoring policy serves the full intent
spectrum, matching a per-intent-strongest heuristic on each objective with a single model. Making that
conditioning effective requires modulation: FiLM re-ranks paths $2.5\times$ more than value-stream
concatenation and wins on the intent that most demands re-ranking, whereas concatenation under-reacts
through common-mode cancellation. The learned selector reaches widest-path goodput while probing roughly
$40\times$ less, reproducing the near-oracle-quality-at-low-probing property it builds on. And even this
near-optimal single-path selector loses about a quarter of its goodput under heavy congestion, because a
single path is a hard capacity bound -- the ceiling that Part~II sets out to break.

\gap{The comparison in \cref{sec:p1eval:overhead,sec:p1eval:ceiling} covers shortest-path, widest-path,
lowest-latency, and random; add the ECMP and SCION-default selectors to complete the six-baseline set of
\cref{sec:p1impl:baselines}, and report a per-intent baseline table once those runs complete.}
